\documentclass[11pt,openany]{book}

% =====================================================
% Fonts & Encoding (CRITICAL FIXES)
% =====================================================
\usepackage[T1]{fontenc}
\usepackage[utf8]{inputenc}
\usepackage{lmodern}                   % Scalable fonts
\usepackage[expansion=false]{microtype}% Avoid pdfTeX font-expansion crash
\usepackage{textcomp}

% =====================================================
% Layout & Typography
% =====================================================
\usepackage[a4paper,margin=1in]{geometry}
\usepackage{setspace}
\onehalfspacing

% =====================================================
% Math & Symbols
% =====================================================
\usepackage{amsmath,amssymb}

% =====================================================
% Figures & Plots
% =====================================================
\usepackage{graphicx}
\usepackage{float}
\usepackage{caption}
\usepackage{subcaption}

% TikZ/PGFPlots (keep available for future real figures)
\usepackage{tikz}
\usepackage{pgfplots}
\pgfplotsset{compat=1.18}

% =====================================================
% Code Listings (UTF-8 SAFE)
% =====================================================
\usepackage{listings}
\newcommand{\SafeListing}[1]{\IfFileExists{#1}{\SafeListing{#1}}{\textit{(Listing not found: #1)}}}
\usepackage{xcolor}
\usepackage{shellesc}

\definecolor{codebg}{rgb}{0.96,0.96,0.96}
\definecolor{codekw}{rgb}{0.15,0.15,0.6}

\lstdefinelanguage{PythonSafe}{
  language=Python,
  morekeywords={dataclass,asdict},
  sensitive=true
}

\lstset{
  backgroundcolor=\color{codebg},
  basicstyle=\ttfamily\small,
  keywordstyle=\color{codekw}\bfseries,
  commentstyle=\itshape\color{gray},
  stringstyle=\color{teal},
  frame=single,
  breaklines=true,
  showstringspaces=false,
  columns=fullflexible,
  keepspaces=true,
  inputencoding=utf8,
  extendedchars=true,
  literate=
    {—}{{---}}1
    {–}{{--}}1
    {’}{{'}}1
    {‘}{{`}}1
    {“}{{``}}1
    {”}{{''}}1
    {…}{{\ldots}}1
}

% =====================================================
% Links & References
% =====================================================
\usepackage{hyperref}
\hypersetup{
  colorlinks=true,
  linkcolor=blue,
  urlcolor=blue,
  citecolor=blue,
  pdfauthor={Jonathan Davanzo},
  pdftitle={End-to-End EEG-Based Prosthetic Finger Control}
}

% =====================================================
% Convenience macros
% =====================================================
\newcommand{\todofig}[2]{%
\begin{figure}[H]
\centering
\fbox{\parbox{0.9\textwidth}{\vspace{0.6em}\textbf{Figure placeholder:} #1\\[0.25em]
\textbf{Insert file:} \texttt{\detokenize{#2}}\\[0.25em]
\vspace{0.6em}}}
\caption{#1}
\end{figure}
}

\newcommand{\exercise}[2]{%
\section*{Exercise #1}
#2
}

\newcommand{\solution}[1]{%
\subsection*{Solution}
#1
}

% =====================================================
% Title
% =====================================================
\title{
\textbf{End-to-End EEG-Based Prosthetic Finger Control}\\
\vspace{0.5em}
\large From Neural Signals to Real-Time Actuation
}
\author{Jonathan Davanzo}
\date{\today}
\ShellEscape{mkdir -p exercises exercises/out figs}

\begin{document}

% =====================================================
% Generated exercise scripts (written at compile time)
% =====================================================
\begin{filecontents*}{exercises/ex00_run_all_sanity.py}
#!/usr/bin/env python3
import subprocess
import sys
from pathlib import Path


def main():
    root = Path(__file__).resolve().parent
    scripts = [
        "ex03_artifact_probe.py",
        "ex04_timebase_break_and_measure.py",
        "ex05_guard_band_sweep.py",
        "ex06_gap_policy_sim.py",
        "ex08_cache_validator_extension.py",
        "ex09_stability_gate_demo.py",
        "ex10_lsl_sanity_probe.py",
        "ex11_backwards_debounce_sim.py",
        "ex12_stream_selector_demo.py",
    ]
    failures = 0
    for script in scripts:
        path = root / script
        result = subprocess.run(
            [sys.executable, str(path), "--synthetic"],
            check=False,
        )
        if result.returncode != 0:
            failures += 1
            print(f"FAILED: {script}")
        else:
            print(f"OK: {script}")
    return 1 if failures else 0


if __name__ == "__main__":
    raise SystemExit(main())
\end{filecontents*}

\begin{filecontents*}{exercises/ex03_artifact_probe.py}
#!/usr/bin/env python3
import argparse
import csv
import json
import math
import random
from pathlib import Path


def load_csv(path):
    rows = []
    with open(path, newline="") as handle:
        reader = csv.reader(handle)
        for row in reader:
            if not row:
                continue
            try:
                values = [float(x) for x in row]
            except ValueError:
                continue
            rows.append(values)
    if not rows:
        raise ValueError("No numeric rows found in CSV.")
    width = min(len(r) for r in rows)
    return [r[:width] for r in rows]


def synthesize(channels, samples, seed):
    random.seed(seed)
    data = []
    for _ in range(samples):
        data.append([random.gauss(0.0, 1.0) for _ in range(channels)])
    start = int(0.4 * samples)
    end = int(0.5 * samples)
    for i in range(start, end):
        for c in range(channels):
            data[i][c] += 8.0
    return data


def sliding_rms_flags(data, fs, window_s, step_s, threshold):
    window = max(1, int(round(window_s * fs)))
    step = max(1, int(round(step_s * fs)))
    flags = []
    for start in range(0, len(data) - window + 1, step):
        end = start + window
        max_rms = 0.0
        for ch in range(len(data[0])):
            acc = 0.0
            for idx in range(start, end):
                value = data[idx][ch]
                acc += value * value
            rms = math.sqrt(acc / window)
            if rms > max_rms:
                max_rms = rms
        flags.append(
            {
                "start_s": start / fs,
                "end_s": end / fs,
                "rms": max_rms,
                "flag": max_rms >= threshold,
            }
        )
    return flags


def main():
    parser = argparse.ArgumentParser(
        description="Sliding RMS artifact probe with a simple threshold."
    )
    parser.add_argument("--csv", help="CSV with one row per sample and channels as columns.")
    parser.add_argument("--fs", type=float, default=256.0, help="Sample rate in Hz.")
    parser.add_argument("--window-s", type=float, default=0.25, help="RMS window in seconds.")
    parser.add_argument("--step-s", type=float, default=0.05, help="Step in seconds.")
    parser.add_argument("--threshold", type=float, default=3.0, help="RMS threshold.")
    parser.add_argument("--channels", type=int, default=4, help="Synthetic channel count.")
    parser.add_argument("--duration-s", type=float, default=10.0, help="Synthetic duration.")
    parser.add_argument("--synthetic", action="store_true", help="Force synthetic data.")
    parser.add_argument("--out-dir", default="exercises/out", help="Output directory.")
    args = parser.parse_args()

    if args.csv and not args.synthetic:
        data = load_csv(args.csv)
    else:
        samples = max(1, int(args.duration_s * args.fs))
        data = synthesize(args.channels, samples, seed=0)

    flags = sliding_rms_flags(
        data=data,
        fs=args.fs,
        window_s=args.window_s,
        step_s=args.step_s,
        threshold=args.threshold,
    )
    flagged = [f for f in flags if f["flag"]]

    out_dir = Path(args.out_dir)
    out_dir.mkdir(parents=True, exist_ok=True)
    out_path = out_dir / "ex03_artifact_probe_intervals.json"
    with out_path.open("w", encoding="utf-8") as handle:
        json.dump({"flags": flags, "flagged_count": len(flagged)}, handle, indent=2)

    print(f"windows={len(flags)} flagged={len(flagged)} out={out_path}")
    return 0


if __name__ == "__main__":
    raise SystemExit(main())
\end{filecontents*}

\begin{filecontents*}{exercises/ex04_timebase_break_and_measure.py}
#!/usr/bin/env python3
import argparse
import json
from pathlib import Path


def build_events(duration_s, event_count, event_len_s):
    if event_count <= 0:
        return []
    gap = max(0.0, (duration_s - event_count * event_len_s) / (event_count + 1))
    events = []
    t = gap
    for _ in range(event_count):
        events.append((t, t + event_len_s))
        t += event_len_s + gap
    return events


def overlap_s(a0, a1, b0, b1):
    return max(0.0, min(a1, b1) - max(a0, b0))


def window_kept_indices(events, duration_s, window_s, step_s, min_overlap_s):
    total = 0
    kept = set()
    start = 0.0
    index = 0
    while start + window_s <= duration_s + 1e-9:
        total += 1
        end = start + window_s
        for e0, e1 in events:
            if overlap_s(start, end, e0, e1) >= min_overlap_s:
                kept.add(index)
                break
        start += step_s
        index += 1
    return total, kept


def main():
    parser = argparse.ArgumentParser(
        description="Simulate timebase misalignment and measure overlap drops."
    )
    parser.add_argument("--duration-s", type=float, default=20.0)
    parser.add_argument("--event-count", type=int, default=5)
    parser.add_argument("--event-len-s", type=float, default=1.0)
    parser.add_argument("--window-s", type=float, default=1.0)
    parser.add_argument("--step-s", type=float, default=0.2)
    parser.add_argument("--min-overlap-frac", type=float, default=0.25)
    parser.add_argument("--offset-s", type=float, default=0.2)
    parser.add_argument("--synthetic", action="store_true", help="Synthetic-only simulation.")
    parser.add_argument("--out-dir", default="exercises/out")
    args = parser.parse_args()

    events = build_events(args.duration_s, args.event_count, args.event_len_s)
    min_overlap_s = args.min_overlap_frac * args.window_s

    aligned_total, aligned_set = window_kept_indices(
        events, args.duration_s, args.window_s, args.step_s, min_overlap_s
    )
    offset_events = [(s + args.offset_s, e + args.offset_s) for s, e in events]
    offset_total, offset_set = window_kept_indices(
        offset_events, args.duration_s, args.window_s, args.step_s, min_overlap_s
    )

    aligned_kept = len(aligned_set)
    offset_kept = len(offset_set)
    aligned_only = len(aligned_set - offset_set)
    offset_only = len(offset_set - aligned_set)
    summary = {
        "aligned_kept": aligned_kept,
        "offset_kept": offset_kept,
        "total_windows": aligned_total,
        "aligned_only": aligned_only,
        "offset_only": offset_only,
        "offset_s": args.offset_s,
    }

    out_dir = Path(args.out_dir)
    out_dir.mkdir(parents=True, exist_ok=True)
    out_path = out_dir / "ex04_timebase_break_and_measure.json"
    with out_path.open("w", encoding="utf-8") as handle:
        json.dump(summary, handle, indent=2)

    print(
        "aligned_kept={aligned_kept} offset_kept={offset_kept} "
        "aligned_only={aligned_only} offset_only={offset_only} out={out}".format(
            out=out_path, **summary
        )
    )
    return 0


if __name__ == "__main__":
    raise SystemExit(main())
\end{filecontents*}

\begin{filecontents*}{exercises/ex05_guard_band_sweep.py}
#!/usr/bin/env python3
import argparse
import json
from pathlib import Path


def parse_list(text):
    return [float(x.strip()) for x in text.split(",") if x.strip()]


def build_events(duration_s, event_count, event_len_s):
    if event_count <= 0:
        return []
    gap = max(0.0, (duration_s - event_count * event_len_s) / (event_count + 1))
    events = []
    t = gap
    for _ in range(event_count):
        events.append((t, t + event_len_s))
        t += event_len_s + gap
    return events


def overlap_s(a0, a1, b0, b1):
    return max(0.0, min(a1, b1) - max(a0, b0))


def sweep(events, duration_s, window_s, step_s, min_overlap_s, guard_s):
    counts = {"total": 0, "kept": 0, "drop_no_overlap": 0, "drop_guard": 0}
    start = 0.0
    while start + window_s <= duration_s + 1e-9:
        counts["total"] += 1
        end = start + window_s
        kept = False
        guard_drop = False
        for e0, e1 in events:
            if overlap_s(start, end, e0, e1) >= min_overlap_s:
                if guard_s > 0.0 and (start < e0 + guard_s or end > e1 - guard_s):
                    guard_drop = True
                else:
                    kept = True
                    break
        if kept:
            counts["kept"] += 1
        elif guard_drop:
            counts["drop_guard"] += 1
        else:
            counts["drop_no_overlap"] += 1
        start += step_s
    return counts


def main():
    parser = argparse.ArgumentParser(
        description="Sweep guard bands and minimum overlap on synthetic events."
    )
    parser.add_argument("--duration-s", type=float, default=20.0)
    parser.add_argument("--event-count", type=int, default=5)
    parser.add_argument("--event-len-s", type=float, default=1.0)
    parser.add_argument("--window-s", type=float, default=1.0)
    parser.add_argument("--step-s", type=float, default=0.2)
    parser.add_argument("--guard-s", default="0.0,0.05,0.1")
    parser.add_argument("--min-overlap-frac", default="0.1,0.25,0.4")
    parser.add_argument("--synthetic", action="store_true", help="Synthetic-only simulation.")
    parser.add_argument("--out-dir", default="exercises/out")
    args = parser.parse_args()

    guard_values = parse_list(args.guard_s)
    overlap_values = parse_list(args.min_overlap_frac)
    events = build_events(args.duration_s, args.event_count, args.event_len_s)

    results = []
    for guard_s in guard_values:
        for overlap in overlap_values:
            counts = sweep(
                events,
                args.duration_s,
                args.window_s,
                args.step_s,
                overlap * args.window_s,
                guard_s,
            )
            entry = {"guard_s": guard_s, "min_overlap_frac": overlap, **counts}
            results.append(entry)
            print(
                "guard_s={guard_s:.3f} min_overlap={min_overlap_frac:.2f} "
                "kept={kept} drop_guard={drop_guard} drop_no_overlap={drop_no_overlap}".format(
                    **entry
                )
            )

    out_dir = Path(args.out_dir)
    out_dir.mkdir(parents=True, exist_ok=True)
    out_path = out_dir / "ex05_guard_band_sweep.json"
    with out_path.open("w", encoding="utf-8") as handle:
        json.dump({"results": results}, handle, indent=2)

    return 0


if __name__ == "__main__":
    raise SystemExit(main())
\end{filecontents*}

\begin{filecontents*}{exercises/ex06_gap_policy_sim.py}
#!/usr/bin/env python3
import argparse
import json
from pathlib import Path


def build_times(duration_s, fs, gap_starts, gap_s):
    dt = 1.0 / fs
    count = max(1, int(duration_s * fs))
    times = [i * dt for i in range(count)]
    for gap_start in sorted(gap_starts):
        for i, t in enumerate(times):
            if t >= gap_start:
                times[i] = t + gap_s
    return times


def window_gap_flags(times, fs, window_s, step_s, gap_threshold_s):
    window = max(2, int(round(window_s * fs)))
    step = max(1, int(round(step_s * fs)))
    flags = []
    for start in range(0, len(times) - window + 1, step):
        window_times = times[start : start + window]
        max_dt = max(
            b - a for a, b in zip(window_times[:-1], window_times[1:])
        )
        flags.append({"index": start, "gap": max_dt > gap_threshold_s})
    return flags


def main():
    parser = argparse.ArgumentParser(
        description="Simulate gaps in timestamps and apply a gap policy."
    )
    parser.add_argument("--duration-s", type=float, default=20.0)
    parser.add_argument("--fs", type=float, default=128.0)
    parser.add_argument("--window-s", type=float, default=1.0)
    parser.add_argument("--step-s", type=float, default=0.25)
    parser.add_argument("--gap-s", type=float, default=0.3)
    parser.add_argument("--gap-threshold-s", type=float, default=0.1)
    parser.add_argument("--gap-starts", default="5,12")
    parser.add_argument("--synthetic", action="store_true", help="Synthetic-only simulation.")
    parser.add_argument("--out-dir", default="exercises/out")
    args = parser.parse_args()

    gap_starts = [float(x.strip()) for x in args.gap_starts.split(",") if x.strip()]
    times = build_times(args.duration_s, args.fs, gap_starts, args.gap_s)
    flags = window_gap_flags(
        times, args.fs, args.window_s, args.step_s, args.gap_threshold_s
    )

    total = len(flags)
    gap_count = sum(1 for f in flags if f["gap"])
    summary = {
        "total_windows": total,
        "gap_windows": gap_count,
        "gap_threshold_s": args.gap_threshold_s,
    }

    out_dir = Path(args.out_dir)
    out_dir.mkdir(parents=True, exist_ok=True)
    out_path = out_dir / "ex06_gap_policy_sim.json"
    with out_path.open("w", encoding="utf-8") as handle:
        json.dump({"summary": summary, "flags": flags}, handle, indent=2)

    print(f"total_windows={total} gap_windows={gap_count} out={out_path}")
    return 0


if __name__ == "__main__":
    raise SystemExit(main())
\end{filecontents*}

\begin{filecontents*}{exercises/ex08_cache_validator_extension.py}
#!/usr/bin/env python3
import argparse
import json
from pathlib import Path


def load_dataset_info_from_npz(path):
    try:
        import numpy as np
    except ImportError as exc:
        raise RuntimeError("numpy is required to load .npz files") from exc

    with np.load(path, allow_pickle=True) as data:
        if "dataset_info" not in data.files:
            raise ValueError("dataset_info not found in npz")
        raw = data["dataset_info"]
        if hasattr(raw, "tolist"):
            raw = raw.tolist()
        if isinstance(raw, list):
            raw = raw[0] if raw else ""
        if isinstance(raw, (bytes, bytearray)):
            raw = raw.decode("utf-8")
        if not isinstance(raw, str):
            raw = str(raw)
        return json.loads(raw)


def load_json(path):
    with open(path, "r", encoding="utf-8") as handle:
        return json.load(handle)


def extra_checks(cache_info, current_info):
    reasons = []
    for key in ("timebase_version", "target_fs"):
        if key not in cache_info or key not in current_info:
            reasons.append(f"missing_{key}")
        elif str(cache_info.get(key)) != str(current_info.get(key)):
            reasons.append(f"{key}_mismatch")
    return reasons


def synthetic_infos(mismatch):
    cache = {"timebase_version": "v1", "target_fs": 256.0, "note": "cache"}
    current = {"timebase_version": "v1", "target_fs": 256.0, "note": "current"}
    if mismatch:
        current["target_fs"] = 128.0
    return cache, current


def main():
    parser = argparse.ArgumentParser(
        description="Validate cached predictions with an extra dataset_info check."
    )
    parser.add_argument("--pred-npz", help="Path to cached prediction npz")
    parser.add_argument("--dataset-info", help="JSON file for current dataset_info")
    parser.add_argument("--synthetic", action="store_true", help="Force synthetic mode")
    parser.add_argument("--mismatch", action="store_true", help="Force a synthetic mismatch")
    parser.add_argument("--out-dir", default="exercises/out")
    args = parser.parse_args()

    cache_info = None
    current_info = None
    if args.pred_npz and not args.synthetic:
        cache_info = load_dataset_info_from_npz(args.pred_npz)
    if args.dataset_info and not args.synthetic:
        current_info = load_json(args.dataset_info)

    if cache_info is None or current_info is None:
        cache_info, current_info = synthetic_infos(args.mismatch)

    reasons = extra_checks(cache_info, current_info)
    status = "pass" if not reasons else "fail"

    out_dir = Path(args.out_dir)
    out_dir.mkdir(parents=True, exist_ok=True)
    out_path = out_dir / "ex08_cache_validator_extension.json"
    with out_path.open("w", encoding="utf-8") as handle:
        json.dump(
            {
                "status": status,
                "reasons": reasons,
                "cache_info": cache_info,
                "current_info": current_info,
            },
            handle,
            indent=2,
        )

    print(f"status={status} reasons={reasons} out={out_path}")
    return 0 if status == "pass" else 2


if __name__ == "__main__":
    raise SystemExit(main())
\end{filecontents*}

\begin{filecontents*}{exercises/ex09_stability_gate_demo.py}
#!/usr/bin/env python3
import argparse
import json
import random
from pathlib import Path


def random_probs(classes):
    values = [random.random() for _ in range(classes)]
    total = sum(values)
    return [v / total for v in values]


def main():
    parser = argparse.ArgumentParser(
        description="Simulate probabilities and apply a stability gate."
    )
    parser.add_argument("--frames", type=int, default=50)
    parser.add_argument("--classes", type=int, default=3)
    parser.add_argument("--confidence", type=float, default=0.7)
    parser.add_argument("--stability-frames", type=int, default=3)
    parser.add_argument("--seed", type=int, default=0)
    parser.add_argument("--synthetic", action="store_true", help="Synthetic-only simulation.")
    parser.add_argument("--out-dir", default="exercises/out")
    args = parser.parse_args()

    random.seed(args.seed)
    history = []
    allowed = []

    for frame in range(args.frames):
        probs = random_probs(args.classes)
        top = max(range(args.classes), key=lambda i: probs[i])
        conf = probs[top]
        history.append(top)
        if len(history) > args.stability_frames:
            history.pop(0)
        stability_ok = (
            len(history) == args.stability_frames and len(set(history)) == 1
        )
        allow = stability_ok and conf >= args.confidence
        allowed.append({"frame": frame, "action": top, "confidence": conf, "allow": allow})

    allow_count = sum(1 for entry in allowed if entry["allow"])
    summary = {
        "frames": args.frames,
        "allow_count": allow_count,
        "confidence": args.confidence,
        "stability_frames": args.stability_frames,
    }

    out_dir = Path(args.out_dir)
    out_dir.mkdir(parents=True, exist_ok=True)
    out_path = out_dir / "ex09_stability_gate_demo.json"
    with out_path.open("w", encoding="utf-8") as handle:
        json.dump({"summary": summary, "events": allowed}, handle, indent=2)

    print(f"allow_count={allow_count} out={out_path}")
    return 0


if __name__ == "__main__":
    raise SystemExit(main())
\end{filecontents*}

\begin{filecontents*}{exercises/ex10_lsl_sanity_probe.py}
#!/usr/bin/env python3
import argparse
import time
from dataclasses import dataclass
from typing import Optional


@dataclass
class ProbeStats:
    samples: int = 0
    backwards_soft: int = 0
    backwards_hard: int = 0
    max_backwards_delta_s: float = 0.0
    gaps: int = 0
    max_gap_s: float = 0.0
    median_dt_s: Optional[float] = None


def median(values):
    values = sorted(values)
    if not values:
        return None
    mid = len(values) // 2
    if len(values) % 2 == 1:
        return float(values[mid])
    return 0.5 * (float(values[mid - 1]) + float(values[mid]))


def main():
    parser = argparse.ArgumentParser(description="Probe an LSL stream for timing issues.")
    parser.add_argument("--seconds", type=float, default=15.0)
    parser.add_argument("--epsilon-s", type=float, default=0.010, help="Soft backwards band.")
    parser.add_argument("--hard-backwards-s", type=float, default=0.200, help="Hard backwards threshold.")
    parser.add_argument("--gap-threshold-s", type=float, default=0.250, help="Gap threshold.")
    parser.add_argument("--synthetic", action="store_true", help="No pylsl required.")
    args = parser.parse_args()

    stats = ProbeStats()
    dts = []
    prev = None

    if args.synthetic:
        ts_series = []
        x = 0.0
        for _ in range(int(args.seconds * 256)):
            x += 1.0 / 256.0
            ts_series.append(x)
        if len(ts_series) > 200:
            ts_series[100] = ts_series[99] - 0.005  # soft backwards
            ts_series[160] = ts_series[159] - 0.6  # hard backwards

        for ts in ts_series:
            if prev is not None:
                dt = ts - prev
                dts.append(dt)
                if dt > args.gap_threshold_s:
                    stats.gaps += 1
                    stats.max_gap_s = max(stats.max_gap_s, dt)
                if ts < prev:
                    delta = prev - ts
                    stats.max_backwards_delta_s = max(stats.max_backwards_delta_s, delta)
                    if delta >= args.hard_backwards_s:
                        stats.backwards_hard += 1
                    elif delta <= args.epsilon_s:
                        stats.backwards_soft += 1
            prev = ts
            stats.samples += 1
    else:
        try:
            from pylsl import StreamInlet, resolve_streams
        except Exception as exc:
            raise SystemExit("pylsl not available. Install pylsl or run with --synthetic.") from exc

        streams = resolve_streams()
        if not streams:
            raise SystemExit("No LSL streams found.")

        chosen = None
        for s in streams:
            name = (s.name() or "").lower()
            typ = (s.type() or "").lower()
            if "eeg" in name or typ == "eeg":
                chosen = s
                break
        if chosen is None:
            chosen = streams[0]

        inlet = StreamInlet(chosen)

        drain_end = time.time() + 0.75
        while time.time() < drain_end:
            inlet.pull_sample(timeout=0.0)

        end_time = time.time() + args.seconds
        while time.time() < end_time:
            _sample, ts = inlet.pull_sample(timeout=0.5)
            if ts is None:
                continue
            if prev is not None:
                dt = ts - prev
                dts.append(dt)
                if dt > args.gap_threshold_s:
                    stats.gaps += 1
                    stats.max_gap_s = max(stats.max_gap_s, dt)
                if ts < prev:
                    delta = prev - ts
                    stats.max_backwards_delta_s = max(stats.max_backwards_delta_s, delta)
                    if delta >= args.hard_backwards_s:
                        stats.backwards_hard += 1
                    elif delta <= args.epsilon_s:
                        stats.backwards_soft += 1
            prev = ts
            stats.samples += 1

    stats.median_dt_s = median([dt for dt in dts if dt > 0])
    approx_fs = (1.0 / stats.median_dt_s) if stats.median_dt_s else None

    print("---- LSL SANITY PROBE ----")
    print(f"samples={stats.samples} median_dt_s={stats.median_dt_s} approx_fs={approx_fs}")
    print(
        f"backwards_soft={stats.backwards_soft} backwards_hard={stats.backwards_hard} "
        f"max_backwards_delta_s={stats.max_backwards_delta_s}"
    )
    print(f"gaps={stats.gaps} max_gap_s={stats.max_gap_s}")
    return 0


if __name__ == "__main__":
    raise SystemExit(main())
\end{filecontents*}

\begin{filecontents*}{exercises/ex11_backwards_debounce_sim.py}
#!/usr/bin/env python3
import argparse
import random
import time
from collections import deque


def clamp_result(prev_mono, raw, hard_backwards_s):
    if prev_mono is None:
        return raw, False, 0.0, False
    if raw >= prev_mono:
        return raw, False, 0.0, False
    delta = prev_mono - raw
    hard = delta >= hard_backwards_s
    return prev_mono, True, delta, hard


def should_break(backwards_events, now_mono, *, hard_backwards, soft_limit, window_s):
    if hard_backwards:
        return True
    backwards_events.append(now_mono)
    while backwards_events and (now_mono - backwards_events[0]) > window_s:
        backwards_events.popleft()
    return len(backwards_events) >= soft_limit


def main():
    parser = argparse.ArgumentParser(description="Simulate debounced backwards policy.")
    parser.add_argument("--steps", type=int, default=500)
    parser.add_argument("--fs", type=float, default=256.0)
    parser.add_argument("--hard-backwards-s", type=float, default=0.200)
    parser.add_argument("--soft-limit", type=int, default=6)
    parser.add_argument("--window-s", type=float, default=1.0)
    parser.add_argument("--p-soft", type=float, default=0.03)
    parser.add_argument("--p-hard", type=float, default=0.002)
    parser.add_argument("--seed", type=int, default=0)
    parser.add_argument("--synthetic", action="store_true")
    args = parser.parse_args()

    random.seed(args.seed)

    raw = 0.0
    prev_mono = None
    backwards_events = deque(maxlen=256)
    segment = 0

    soft_count = 0
    hard_count = 0
    breaks = 0

    for i in range(args.steps):
        dt = 1.0 / args.fs
        raw += dt

        r = random.random()
        if r < args.p_hard:
            raw -= 0.6
        elif r < args.p_hard + args.p_soft:
            raw -= 0.005

        mono, clamped, delta, hard = clamp_result(prev_mono, raw, args.hard_backwards_s)
        if clamped:
            if hard:
                hard_count += 1
            else:
                soft_count += 1

        now_mono = time.monotonic()
        do_break = should_break(
            backwards_events,
            now_mono,
            hard_backwards=hard,
            soft_limit=args.soft_limit,
            window_s=args.window_s,
        )
        if do_break:
            breaks += 1
            segment += 1
            backwards_events.clear()
            prev_mono = None
            print(f"[break] step={i} reason={'hard' if hard else 'burst'} delta={delta:.3f} seg={segment}")
            continue

        prev_mono = mono

    print("---- DEBOUNCE SUMMARY ----")
    print(f"soft_clamps={soft_count} hard_clamps={hard_count} segment_breaks={breaks}")
    return 0


if __name__ == "__main__":
    raise SystemExit(main())
\end{filecontents*}

\begin{filecontents*}{exercises/ex12_stream_selector_demo.py}
#!/usr/bin/env python3
import argparse


def pick_candidate(candidates, *, name_contains=None, type_equals=None, min_channels=4, require_unique=True):
    filtered = []
    for c in candidates:
        if type_equals is not None and str(c.get("type", "")).lower() != str(type_equals).lower():
            continue
        if name_contains is not None and str(name_contains).lower() not in str(c.get("name", "")).lower():
            continue
        if int(c.get("channels", 0)) < int(min_channels):
            continue
        filtered.append(c)

    if not filtered:
        raise RuntimeError(f"No streams match selector. All candidates={candidates}")

    if require_unique and len(filtered) != 1:
        raise RuntimeError(
            "Ambiguous stream selection. Candidates matched selector:\n"
            + "\n".join([f"  - {x}" for x in filtered])
            + "\nFix by providing a stricter selector (name_contains or type_equals)."
        )
    return filtered[0]


def main():
    parser = argparse.ArgumentParser(description="Pure stream selection demo.")
    parser.add_argument("--synthetic", action="store_true")
    _ = parser.parse_args()

    candidates = [
        {"name": "MuseS-EEG", "type": "EEG", "channels": 5, "uid": "AAA"},
        {"name": "SimEEG", "type": "EEG", "channels": 8, "uid": "BBB"},
        {"name": "MarkerStream", "type": "Markers", "channels": 1, "uid": "CCC"},
    ]

    print("All candidates:")
    for c in candidates:
        print(" ", c)

    try:
        chosen = pick_candidate(
            candidates,
            type_equals="EEG",
            min_channels=4,
            require_unique=True,
        )
        print("chosen:", chosen)
    except RuntimeError as e:
        print("EXPECTED ambiguity error:")
        print(e)

    chosen = pick_candidate(
        candidates,
        name_contains="Muse",
        type_equals="EEG",
        min_channels=4,
        require_unique=True,
    )
    print("chosen after disambiguation:", chosen)
    return 0


if __name__ == "__main__":
    raise SystemExit(main())
\end{filecontents*}

\begin{filecontents*}{make_figs.py}
#!/usr/bin/env python3
import json
from pathlib import Path

def load_json(path: Path):
    with path.open("r", encoding="utf-8") as f:
        return json.load(f)

def main():
    try:
        import matplotlib.pyplot as plt
    except Exception as e:
        raise SystemExit("matplotlib required for figure generation") from e

    repo_root = Path(__file__).resolve().parent

    # Where LaTeX expects PDFs to live
    figs_dir = repo_root / "figs"
    figs_dir.mkdir(parents=True, exist_ok=True)

    out_dir = repo_root / "exercises" / "out"

    wrote = []

    # --- Fig: timebase effect (ex04) ---
    p = out_dir / "ex04_timebase_break_and_measure.json"
    if p.exists():
        d = load_json(p)
        labels = ["aligned_kept", "offset_kept", "aligned_only", "offset_only"]
        vals = [float(d.get(k, 0) or 0) for k in labels]
        plt.figure()
        plt.bar(labels, vals)
        plt.xticks(rotation=20, ha="right")
        plt.ylabel("count")
        plt.title(f"Timebase misalignment impact (offset={d.get('offset_s')})")
        plt.tight_layout()
        out_path = figs_dir / "timebase_alignment_effect.pdf"
        plt.savefig(out_path)
        plt.close()
        wrote.append(out_path)

    # --- Fig: stability gate (ex09) ---
    p = out_dir / "ex09_stability_gate_demo.json"
    if p.exists():
        d = load_json(p)
        events = d.get("events", []) or []
        allow = [1 if (e.get("allow") is True) else 0 for e in events]
        plt.figure()
        plt.plot(allow)
        plt.ylim(-0.1, 1.1)
        plt.xlabel("frame")
        plt.ylabel("allow (0/1)")
        plt.title("Stability gate allow/deny over time")
        plt.tight_layout()
        out_path = figs_dir / "stability_gate_timeline.pdf"
        plt.savefig(out_path)
        plt.close()
        wrote.append(out_path)

    if wrote:
        print("Wrote:")
        for p in wrote:
            print(" ", p)
    else:
        print("No figures written (missing exercises/out JSON outputs).")
    return 0

if __name__ == "__main__":
    raise SystemExit(main())
\end{filecontents*}

\maketitle
\tableofcontents
\clearpage

\chapter*{How to Use This Course}
\addcontentsline{toc}{chapter}{How to Use This Course}

This is a self-paced, upper-undergraduate to graduate course on end-to-end
EEG-based finger control. Each coding exercise has a minimal standalone script
that is emitted when you compile this document. The scripts are written to
\texttt{\detokenize{exercises/}} via \texttt{\detokenize{filecontents*}} blocks in the LaTeX source.

Compile the document (this writes the scripts):
\begin{lstlisting}
pdflatex -interaction=nonstopmode eeg_hand_lessons.tex
\end{lstlisting}

Run the full sanity sweep or a single exercise script:
\begin{lstlisting}
python exercises/ex00_run_all_sanity.py
python exercises/ex03_artifact_probe.py --synthetic
\end{lstlisting}

Generate figures from exercise outputs:
\begin{lstlisting}
python make_figs.py
\end{lstlisting}

By default, each script writes outputs to \texttt{\detokenize{exercises/out/}} and accepts
CLI flags to point at real data when you have it.

% =====================================================
\part{Foundations}
% =====================================================

\chapter{Lesson 1: What Are We Building and Why?}

\section{The Goal}

The objective of this project is to infer \emph{individual finger movements}
from non-invasive EEG signals in real time, with sufficient reliability to
control a physical prosthetic hand.

This is a \textbf{real-time neuroengineering system}, not a static
classification task. Success requires both:
\begin{itemize}
  \item \textbf{signal integrity} (time alignment, label correctness, drift control),
  \item \textbf{decision integrity} (uncertainty-aware inference, safe actuation).
\end{itemize}

\section{System Contract (What Must Always Be True)}

We define a few invariants that the system must maintain:
\begin{itemize}
  \item \textbf{Timebase invariant:} samples and events must be expressed in the same clock domain.
  \item \textbf{Label invariant:} the extracted windows must overlap the intended event with a defined policy.
  \item \textbf{Evaluation invariant:} train/test splits and cached predictions must not mix samples across incompatible datasets.
  \item \textbf{Safety invariant:} actuation is gated by confidence \emph{and} stability over time.
\end{itemize}

\todofig{End-to-end pipeline overview (data acquisition $\rightarrow$ features $\rightarrow$ windows $\rightarrow$ train $\rightarrow$ eval $\rightarrow$ realtime).}{figs/pipeline_overview.pdf}


\section{Project pipeline: streaming, labeling, windowing, and training}

This project is organized around three scripts that must interoperate correctly:

\begin{enumerate}
  \item \textbf{Step 1 --- \texttt{1\_stream\_and\_record.py}:} Connects to a Muse 2 over BLE, publishes an EEG stream over Lab Streaming Layer (LSL), and simultaneously records two CSV files:
  \begin{itemize}
    \item \texttt{*\_raw.csv}: the continuous EEG time series (4 channels; only the first four channels are used for modeling).
    \item \texttt{*\_events.csv}: event markers (labels) with timestamps in the same timebase as the EEG samples.
  \end{itemize}
  \item \textbf{Step 1b --- \texttt{1b\_extract\_windows.py}:} Loads \texttt{*\_raw.csv} and \texttt{*\_events.csv}, resamples onto a fixed time grid, and extracts fixed-length windows aligned to each event. Any event/window that would contain NaNs (or non-finite values) is dropped to guarantee safe downstream training.
  \item \textbf{Step 2 --- \texttt{2\_train\_model.py}:} Trains the PyTorch model on the extracted windows. This step assumes \emph{no NaNs} are present in the input tensors; therefore Step 1b is responsible for filtering/dropping problematic windows.
\end{enumerate}

\subsection{Why LSL and why timestamps matter}
LSL provides a common clock domain for synchronizing streamed biosignals and event markers across processes and machines \cite{LSL, pylsl}. In this pipeline:

\begin{itemize}
  \item Each EEG sample is associated with an LSL timestamp.
  \item Event times are stored using the same LSL timebase so that Step 1b can align events and EEG samples precisely.
\end{itemize}

\subsection{Recorded data formats}

\paragraph{Raw EEG CSV (\texttt{*\_raw.csv}).}
Step 1 writes a CSV with at least the following columns:
\begin{itemize}
  \item \texttt{time\_s}: session-relative time in seconds (derived from the LSL sample timestamps).
  \item \texttt{lsl\_ts}: the absolute LSL timestamp for each sample.
  \item \texttt{local\_ts}: the local clock timestamp for each sample (for debugging/health checks).
  \item Four EEG channel columns. With Muse 2, the canonical four EEG channels are typically labeled \texttt{TP9}, \texttt{AF7}, \texttt{AF8}, \texttt{TP10} and the sample rate is 256\,Hz \cite{Muse2spec}.
\end{itemize}

\paragraph{Events CSV (\texttt{*\_events.csv}).}
Step 1 writes a CSV with an \texttt{event\_time\_s} column and a \texttt{label} column (and may include additional metadata such as an \texttt{action} or \texttt{note}). Step 1b loads these rows and aligns each event to a time window in the raw EEG.

\subsection{NaN safety invariant}
PyTorch training in Step 2 must never receive NaNs. The project enforces the invariant:

\begin{quote}
\textbf{No NaNs enter Step 2.} If Step 1b encounters NaNs (or non-finite values) in EEG samples, resampling, labels, or window assembly, the affected event/window is dropped before serialization.
\end{quote}

\subsection{End-to-end checklist}
Before collecting important data, verify:

\begin{enumerate}
  \item Step 1 connects to the Muse 2 reliably via BLE (Python BLE support is typically provided by \texttt{bleak} \cite{Bleak}).
  \item Step 1 is writing \texttt{*\_raw.csv} and \texttt{*\_events.csv} continuously during the session (not only at shutdown).
  \item Step 1b resolves input files from \texttt{meta.json} (preferred) or by globbing \texttt{*\_raw.csv} and \texttt{*\_events.csv} in the session directory.
  \item Step 1b drops any NaN/non-finite windows and reports counts.
  \item Step 2 starts training without NaN-related crashes.
\end{enumerate}



\section{Exercises}

\exercise{1 (Conceptual)}{
Why might a model report high accuracy yet be useless for real-time prosthetic control?
List at least three \emph{concrete} failure mechanisms tied to the system contract above.
}

% =====================================================

\chapter{Lesson 2: Neurophysiology for Engineers}

\section{Motor Control: What EEG Can Realistically Encode}

Voluntary finger movement is primarily initiated in the \textbf{primary motor cortex (M1)}.
The cortical representation is somatotopic, but fingers overlap.

EEG is limited by volume conduction: each electrode measures a weighted sum of many sources.
So separability comes from \textbf{temporal structure} (planning $\rightarrow$ execution),
not purely spatial structure.

\todofig{Somatotopy schematic and electrode montage for Muse-class devices.}{figs/montage_somatotopy.pdf}

\section{Motor Rhythms and Movement-Related Signals}

Two commonly exploited phenomena:
\begin{itemize}
\item \textbf{Mu rhythm} (roughly 8--13 Hz): tends to desynchronize around movement planning/execution.
\item \textbf{Beta rhythm} (roughly 13--30 Hz): also modulates with movement and post-movement rebound.
\end{itemize}

In consumer EEG, you rarely get clean cortical rhythms. The system must be robust to contamination.

\section{Exercises}

\exercise{2 (Conceptual)}{
Explain why finger classification is more plausible in a \emph{time window around movement onset}
than at a single instant.
}

% =====================================================
\part{Signals, Time, and Labels}
% =====================================================

\chapter{Lesson 3: EEG Signal Acquisition and Digitization}

\section{Sampling}

Let $x(t)$ be a continuous EEG signal sampled at frequency $f_s$:
\[
x[n] = x\left(\frac{n}{f_s}\right).
\]
The Nyquist limit is $f_s/2$, but practical issues (noise, drift, artifacts) dominate.

\section{Artifacts That Matter in This Project}

You should assume artifacts are always present:
\begin{itemize}
  \item eye blinks / saccades,
  \item facial muscle (EMG),
  \item cable/electrode motion,
  \item posture changes.
\end{itemize}

\todofig{Example raw signal with blink artifact and movement onset overlay.}{figs/raw_with_events.pdf}

\section{Exercises}

\exercise{3 (Coding)}{
Add a quick artifact probe: compute per-channel RMS in a sliding window and log
windows when RMS exceeds a threshold. Use this as a candidate "artifact flag"
for window filtering.

Run: \texttt{\detokenize{python exercises/ex03_artifact_probe.py --synthetic}}
}

% ============================================================

\chapter{Lesson 4: Timebases and Clock Alignment (Critical)}

\section{The Core Insight}

\textbf{Labels are meaningless unless they share the same timebase as the signal.}

In this project, there are two clocks:
\begin{itemize}
  \item \textbf{LSL stream clock} for EEG samples,
  \item \textbf{local system clock} for UI/keyboard events.
\end{itemize}
We must map events into the LSL domain (or map samples into the local domain), then operate in one unified time coordinate.

\section{Stream-Relative Timebase}

Define stream-relative seconds:
\[
t_{\mathrm{stream}} = t_{\mathrm{LSL}} - t_{\mathrm{LSL\,start}}.
\]

At session start, measure both clocks and compute a constant offset:
\[
\Delta = t_{\mathrm{LSL\,start}} - t_{\mathrm{local\,start}}.
\]
Then any local event timestamp $t_{\mathrm{local}}$ maps to the LSL domain:
\[
t_{\mathrm{event\,LSL}} = t_{\mathrm{local}} + \Delta.
\]
And finally to stream-relative seconds:
\[
t_{\mathrm{event\,stream}} = t_{\mathrm{event\,LSL}} - t_{\mathrm{LSL\,start}}.
\]

\todofig{Timebase alignment diagram: local clock and LSL clock mapped by an offset.}{figs/timebase_alignment.pdf}

\section{Why a Constant Offset Is Usually Enough (and When It Isn't)}
If both clocks are stable over a session, a constant offset $\Delta$ maps event times into stream time.
However, drift can occur due to OS scheduling, Bluetooth jitter, or long sessions.
Your \texttt{\detokenize{timebase_selfcheck}} is a runtime statistical test:
it measures nearest-sample deltas for recent events and flags warning/error when alignment degrades.

\section{Engineering Rule}
\textbf{Never train} on data collected under a timebase warning without explicitly labeling it as such.
For real systems, misalignment is not “noise”—it is \emph{systematic label corruption}.

\todofig{Timebase misalignment impact generated from Exercise 4 output.}{figs/timebase_alignment_effect.pdf}

\section{Code Excerpts: Monotonic Clamping and Self-Check}

\begin{lstlisting}[language=PythonSafe,caption={Excerpt from \texttt{\detokenize{utils/timebase.py}} L6--L17}]
def clamp_monotonic_time(
    prev_s: Optional[float],
    current_s: Optional[float],
    epsilon_s: float = 1e-6,
) -> Tuple[Optional[float], bool]:
    if current_s is None:
        return None, False
    if prev_s is None:
        return current_s, False
    if current_s + epsilon_s < prev_s:
        return prev_s, True
    return current_s, False
\end{lstlisting}

\begin{lstlisting}[language=PythonSafe,caption={Excerpt from \texttt{\detokenize{utils/timebase_selfcheck.py}} L19--L65}]
def evaluate_timebase_alignment(
    recent_sample_times: Sequence[float],
    event_times: Iterable[float],
    warn_threshold_s: float = 0.05,
    error_threshold_s: float = 0.2,
) -> TimebaseSelfCheck:
    sample_times = np.asarray(recent_sample_times, dtype=float)
    event_times_arr = np.asarray(list(event_times), dtype=float)
    if sample_times.size == 0 or event_times_arr.size == 0:
        return TimebaseSelfCheck(
            max_abs_delta_s=None,
            mean_abs_delta_s=None,
            warn_threshold_s=warn_threshold_s,
            error_threshold_s=error_threshold_s,
            warn=False,
            error=False,
        )

    deltas = []
    for event_time in event_times_arr:
        nearest = sample_times[np.argmin(np.abs(sample_times - event_time))]
        deltas.append(float(event_time - nearest))
    deltas_arr = np.abs(np.asarray(deltas, dtype=float))
    max_abs = float(np.max(deltas_arr)) if deltas_arr.size else None
    mean_abs = float(np.mean(deltas_arr)) if deltas_arr.size else None
    warn = (
        (max_abs is not None and max_abs >= warn_threshold_s)
        or (mean_abs is not None and mean_abs >= warn_threshold_s)
    )
    error = (
        (max_abs is not None and max_abs >= error_threshold_s)
        or (mean_abs is not None and mean_abs >= error_threshold_s)
        or (
            max_abs is not None
            and mean_abs is not None
            and max_abs >= warn_threshold_s
            and mean_abs >= warn_threshold_s
        )
    )
    return TimebaseSelfCheck(
        max_abs_delta_s=max_abs,
        mean_abs_delta_s=mean_abs,
        warn_threshold_s=warn_threshold_s,
        error_threshold_s=error_threshold_s,
        warn=warn,
        error=error,
    )
\end{lstlisting}

\section{Exercises}

\exercise{4 (Debugging)}{
Simulate timebase misalignment and measure overlap loss. Report:
\begin{itemize}
\item how many windows drop due to no overlap,
\item how the label distribution changes,
\item how model accuracy changes if you train with the misaligned windows.
\end{itemize}

Run: \texttt{\detokenize{python exercises/ex04_timebase_break_and_measure.py --synthetic}}
}

% ============================================================

\chapter{Lesson 5: Event Labeling as a Human--System Interface}

\section{Motor Latency and Fuzzy Labels}

Human movement onset is not a delta function:
\begin{itemize}
\item intention forms (hundreds of ms),
\item planning occurs (hundreds of ms),
\item muscles activate (tens to hundreds of ms),
\item release follows.
\end{itemize}
So we treat events as intervals with uncertainty, not perfect boundaries.

\section{Guard Bands and Overlap Policy}

A simple, defensible extraction policy:
\begin{itemize}
\item include a window if it overlaps an event by at least $\alpha$ of window duration,
\item exclude windows within a transition guard band (avoid ambiguous boundaries),
\item optionally exclude windows overlapping artifact intervals.
\end{itemize}

\todofig{Window overlap geometry: event interval, window interval, overlap fraction, guard bands.}{figs/window_overlap_policy.pdf}

\section{Code Excerpt: Gap and Overlap Policy}

\begin{lstlisting}[language=PythonSafe,caption={Excerpt from \texttt{\detokenize{1b_extract_windows.py}} L1018--L1047}]
        max_dt = (
            float(np.max(np.diff(times_pad))) if times_pad.size >= 2 else float("inf")
        )
        gap_flag = int(max_dt > GAP_THRESHOLD_SEC)
        if gap_flag and not args.allow_gaps:
            drop_gap += 1
            continue

        # Overlaps with events
        overlapping = []
        any_artifact_flag = 0

        for e in events:
            ov = overlap_s(
                window_start, window_end, float(e["onset_s"]), float(e["end_s"])
            )
            if ov <= 0:
                continue

            ov_frac = ov / float(window_sec)

            if e.get("type") == "artifact" and ov_frac >= ARTIFACT_MIN_OVERLAP_FRAC:
                any_artifact_flag = 1
                break

            overlapping.append((ov, ov_frac, e))

        if any_artifact_flag:
            drop_artifact += 1
            continue
\end{lstlisting}

\section{Exercises}

\exercise{5 (Coding)}{
Sweep guard band width and minimum overlap, then observe:
\begin{itemize}
\item total windows kept,
\item class balance,
\item action accuracy vs finger (non-REST) accuracy.
\end{itemize}
Provide one recommended guard band setting and justify it.

Run: \texttt{\detokenize{python exercises/ex05_guard_band_sweep.py --synthetic}}
}

% ============================================================
\part{Models, Training, and Evaluation}
% ============================================================

\chapter{Lesson 6: Window Extraction as Statistical Curation}

\section{Why Windowing Is More Than Preprocessing}

Window extraction is \textbf{dataset design}. It defines what the model is allowed to learn.
Two common pitfalls:
\begin{itemize}
\item \textbf{row-based windows} when sample rate is variable (breaks time meaning),
\item \textbf{label leakage} from future samples if alignment is wrong.
\end{itemize}

\section{Time-based Resampling}

If the logging rate drifts, resample onto a fixed time grid and interpolate.
This preserves fixed-duration windows.

\todofig{Resampling schematic: irregular timestamps $\rightarrow$ uniform grid $\rightarrow$ fixed windows.}{figs/resampling_grid.pdf}

\section{Code Excerpt: Resampling to a Fixed Grid}

\begin{lstlisting}[language=PythonSafe,caption={Excerpt from \texttt{\detokenize{1b_extract_windows.py}} L1150--L1159}]
        if not is_valid_action_finger(action_id, finger_id):
            drop_invalid_label += 1
            continue

        # Resample to fixed window_samples
        grid = np.linspace(window_start, window_end, window_samples, endpoint=False)
        segment = np.empty((window_samples, signal.shape[1]), dtype=float)
        for ch_idx in range(signal.shape[1]):
            segment[:, ch_idx] = np.interp(grid, times_pad, signal_pad[:, ch_idx])
\end{lstlisting}

\section{Exercises}

\exercise{6 (Debugging)}{
Intentionally introduce a gap in timestamps (simulate missing samples).
Show how interpolation and gap policy prevent invalid windows from entering training.

Run: \texttt{\detokenize{python exercises/ex06_gap_policy_sim.py --synthetic}}
}

% ============================================================

\chapter{Lesson 7: Model Architecture and Objectives}

\section{Two-Head Prediction: Action + Finger}

We predict:
\begin{itemize}
  \item \textbf{action} $\in \{\mathrm{REST},\mathrm{OPEN},\mathrm{CLOSE}\}$,
  \item \textbf{finger} $\in \{\mathrm{NONE},\mathrm{THUMB},\dots,\mathrm{PINKY}\}$,
\end{itemize}
often evaluating finger on non-REST samples only.

\section{Why Temporal Models Help}

If per-window inputs contain sequences (e.g., 64 samples $\times$ 4 channels), the model must learn dynamics.
A simple structure: temporal conv layers + recurrent layer (or transformer-lite) feeding two classifiers.

\todofig{Architecture diagram: feature extractor backbone with two output heads.}{figs/model_architecture.pdf}

\section{Exercises}

\exercise{7 (Conceptual)}{
Why can joint training improve finger prediction compared to training finger-only?
}

% ============================================================

\chapter{Lesson X: PyTorch Methods for Neuroengineering Models}

\section{Tensor Shapes Are the Real Interface}
In this project, your model interface is a contract on \emph{shape}, \emph{dtype}, and \emph{time meaning}.
A typical input window is shaped as:
\[
\texttt{x} \in \mathbb{R}^{B \times T \times C},
\]
where $B$ is batch size, $T$ samples per window, and $C$ channels.
If you reorder dimensions (e.g., $B \times C \times T$), do it intentionally and document it.

\section{Dataset and DataLoader}
Use \texttt{\detokenize{torch.utils.data.Dataset}} to define exactly what a “sample” is:
a window tensor + action label + finger label + metadata (subject, session, timebase version).
Then DataLoader handles batching, shuffling, and determinism.

\section{Two-Head Loss: Action + Finger}
Let $f_\theta(x)$ output logits for both heads:
\[
(\ell_f, \ell_a) = f_\theta(x)
\]
where $\ell_f \in \mathbb{R}^{B\times N_f}$ and $\ell_a \in \mathbb{R}^{B\times N_a}$.
Use cross-entropy per head and combine:
\[
\mathcal{L} = \lambda_a \,\mathrm{CE}(\ell_a, y_a) + \lambda_f \,\mathrm{CE}(\ell_f, y_f).
\]
Tune $(\lambda_a,\lambda_f)$ to prevent trivial REST solutions.

\section{Train/Eval Mode Is Not Optional}
\texttt{\detokenize{model.train()}} enables dropout and batchnorm updates.
\texttt{\detokenize{model.eval()}} freezes them.
Incorrect mode usage silently corrupts calibration and uncertainty estimates.

\section{Reproducibility Controls}
At minimum: seed Python, NumPy, and PyTorch; log versions; log dataset hashes/filters.
Reproducibility is part of safety: you must be able to explain a behavior regression.

% ============================================================

\chapter{Lesson 8: Evaluation Without Self-Deception}

\section{Why Accuracy Is a Trap}

Accuracy can look good even if:
\begin{itemize}
  \item the model always predicts REST,
  \item the model is miscalibrated (probabilities meaningless),
  \item cached predictions are mismatched to the current dataset.
\end{itemize}

\section{Decision-Theoretic View of Actuation}
You are not maximizing accuracy; you are minimizing expected harm.
A safe controller triggers only when expected utility is positive:
\begin{itemize}
\item confidence high,
\item uncertainty low,
\item stability satisfied,
\item context constraints satisfied (e.g., not in an artifact interval).
\end{itemize}
Calibration matters because it links reported probability to real-world correctness.

\section{Reliability Diagrams and ECE}
A reliability diagram bins predictions by confidence and plots empirical accuracy per bin.
ECE summarizes the gap between confidence and accuracy across bins.
If you threshold at 0.8 confidence but your 0.8 bin is only 0.55 accurate, your gate is unsafe.

\section{Calibration: Expected Calibration Error (ECE)}

ECE summarizes whether predicted probabilities match empirical correctness.
If you gate actuation on confidence, miscalibration is a safety issue.

\todofig{Reliability diagram: confidence bins vs accuracy (ECE).}{figs/reliability_diagram.pdf}

\section{Cache Validation and Dataset Identity}

Cached prediction reuse is only safe if the cache corresponds to the same dataset identity:
\begin{itemize}
\item same NPZ identity (hash/size/path policy),
\item same filters (subject id, max samples),
\item same test indices (global indices preferred),
\item spot-check labels against current data.
\end{itemize}

\section{Code Excerpts: Cache Validation}

\begin{lstlisting}[language=PythonSafe,caption={Excerpt from \texttt{\detokenize{utils/eval_utils.py}} L54--L113}]
def validate_cached_predictions_with_dataset_info(
    *,
    action_probs: np.ndarray,
    finger_probs: np.ndarray,
    y_action_test: np.ndarray,
    y_finger_test: np.ndarray,
    test_idx: np.ndarray,
    n_actions: int,
    n_fingers: int,
    n_samples_current: int,
    dataset_info_cache: Optional[Dict[str, Any]],
    dataset_info_current: Dict[str, Any],
    y_action_current: np.ndarray,
    y_finger_current: np.ndarray,
    spotcheck_k: int = 10,
    rng_seed: int = 0,
) -> Tuple[bool, List[str]]:
    reasons: List[str] = []
    cache_ok = validate_cached_predictions(
        action_probs=action_probs,
        finger_probs=finger_probs,
        y_action_test=y_action_test,
        y_finger_test=y_finger_test,
        test_idx=test_idx,
        n_actions=n_actions,
        n_fingers=n_fingers,
        n_samples=n_samples_current,
    )
    if not cache_ok:
        reasons.append("cache_shape_invalid")
        return False, reasons

    if dataset_info_cache is None:
        return False, ["legacy_cache_missing_dataset_info"]

    cache_filters = dataset_info_cache.get("filters", {}) or {}
    current_filters = dataset_info_current.get("filters", {}) or {}

    cache_subject_id = (
        ""
        if cache_filters.get("subject_id") is None
        else str(cache_filters.get("subject_id"))
    )
    current_subject_id = (
        ""
        if current_filters.get("subject_id") is None
        else str(current_filters.get("subject_id"))
    )
    if cache_subject_id != current_subject_id:
        reasons.append("filter_subject_id_mismatch")

    cache_max_samples = cache_filters.get("max_samples")
    current_max_samples = current_filters.get("max_samples")
    if cache_max_samples != current_max_samples:
        reasons.append("filter_max_samples_mismatch")

    if dataset_info_cache.get("experiment_hash") != dataset_info_current.get(
        "experiment_hash"
    ):
        reasons.append("experiment_hash_mismatch")
\end{lstlisting}

\begin{lstlisting}[language=PythonSafe,caption={Excerpt from \texttt{\detokenize{utils/eval_utils.py}} L139--L150}]
    if spotcheck_k > 0 and len(test_idx) > 0:
        rng = np.random.default_rng(rng_seed)
        sample_count = min(int(spotcheck_k), len(test_idx))
        positions = rng.choice(len(test_idx), size=sample_count, replace=False)
        for pos in positions:
            idx = int(test_idx[pos])
            if (
                y_action_test[pos] != y_action_current[idx]
                or y_finger_test[pos] != y_finger_current[idx]
            ):
                reasons.append("spotcheck_label_mismatch")
                break
\end{lstlisting}

\begin{lstlisting}[language=PythonSafe,caption={Excerpt from \texttt{\detokenize{3_evaluate_model.py}} L773--L813}]
    cached = _load_predictions_if_present(pred_npz_path)
    cache_rejected_reasons: Optional[List[str]] = None

    cached_used = False
    split_attempts = 0
    if cached is not None:
        action_probs = np.asarray(cached["action_probs"])
        finger_probs = np.asarray(cached["finger_probs"])
        y_action_test = np.asarray(cached["y_action"])
        y_finger_test = np.asarray(cached["y_finger"])
        test_idx = resolve_cached_test_indices(cached)
        if test_idx is None:
            cache_rejected_reasons = ["missing_test_indices"]
            cached = None
        else:
            test_idx = np.asarray(test_idx).astype(np.int64)

        if cached is not None:
            all_idx = np.arange(len(y_action), dtype=np.int64)
            test_mask = np.zeros(len(y_action), dtype=bool)
            test_mask[test_idx] = True
            train_idx = all_idx[~test_mask]

            dataset_info_cache = _parse_dataset_info(cached)
            cache_ok, reject_reasons = validate_cached_predictions_with_dataset_info(
                action_probs=action_probs,
                finger_probs=finger_probs,
                y_action_test=y_action_test,
                y_finger_test=y_finger_test,
                test_idx=test_idx,
                n_actions=n_actions,
                n_fingers=n_fingers,
                n_samples_current=len(y_action),
                dataset_info_cache=dataset_info_cache,
                dataset_info_current=dataset_info_current,
                y_action_current=y_action,
                y_finger_current=y_finger,
                spotcheck_k=10,
                rng_seed=0,
            )

\end{lstlisting}

\section{Exercises}

\exercise{8 (Coding)}{
Add one additional validator check that would catch a subtle mismatch.
Examples:
\begin{itemize}
\item verify class label maps match,
\item verify number of channels and sampling rate match,
\item verify timebase version matches.
\end{itemize}
Implement and explain why it matters.

Run: \texttt{\detokenize{python exercises/ex08_cache_validator_extension.py --synthetic}}
}

% ============================================================

\chapter{Repository and Script Overview}
The public GitHub repository is the canonical source of truth for code structure, configuration, and usage instructions. \cite{repo_github}

\section{Core workflow scripts}
This project’s most critical offline pipeline is:
\begin{enumerate}
\item \textbf{Step 1: Stream and record} (\texttt{1\_stream\_and\_record.py}) ensures a Muse 2 EEG LSL stream is available, and writes:
\begin{itemize}
\item \texttt{\_raw.csv}: raw EEG samples with a consistent timebase, and
\item \texttt{\_events.csv}: event markers in the same time domain.
\end{itemize}
\item \textbf{Step 1b: Extract windows} (\texttt{1b\_extract\_windows.py}) reads the raw/events CSVs, resamples onto a fixed grid, extracts fixed-duration windows aligned to events, and \textbf{drops any windows/events that would contain NaNs} so downstream training is guaranteed NaN-free.
\item \textbf{Step 2: Train model} (\texttt{2\_train\_model.py}) trains the PyTorch classifier on extracted windows. PyTorch training cannot accept NaNs; therefore Step 1b enforces NaN-free tensors before training begins.
\end{enumerate}

\section{Where to start}
If you are new to the codebase, begin with the repository README and then follow the scripts in order (Step 1 $\rightarrow$ Step 1b $\rightarrow$ Step 2). \cite{repo_github}


\part{Real-Time Systems and Safety}
% ============================================================

\chapter{Lesson 9: Real-Time Inference and Stability}

\section{Latency vs Accuracy}

Smaller windows and faster step sizes reduce latency but increase variance.
Stability controls (hysteresis, consecutive-frame agreement) convert noisy per-window outputs into safe decisions.

\todofig{Latency vs accuracy curve with chosen operating point highlighted.}{figs/latency_accuracy_tradeoff.pdf}

\section{Uncertainty (MC Dropout) and Safe Actuation}


Dropout is a regularization technique: during training, we randomly “turn off” a fraction of neurons (or connections) on each forward pass. This reduces overfitting by preventing the network from relying too heavily on any single pathway. In standard inference, dropout is disabled so the model’s output is deterministic.

\textbf{Monte Carlo (MC) Dropout} keeps dropout \emph{enabled during inference} and runs the same input through the network multiple times (e.g., 10--100 forward passes). Each pass uses a different random dropout mask, producing a slightly different prediction. From the resulting set of predictions we compute:
\begin{itemize}
\item a \textbf{final prediction} as the mean (or majority vote) across passes, and
\item an \textbf{uncertainty measure} from the spread (e.g., variance or entropy) of those predictions.
\end{itemize}
When the predictions are tightly clustered, uncertainty is low; when they are spread out, uncertainty is high.

\textbf{High-school level (why it matters).} EEG is noisy, and the user’s physiology and headset contact vary over time. Uncertainty helps us detect when the model “isn’t sure,” so we can avoid acting on unreliable predictions.

\textbf{Undergraduate level (ensemble intuition).} MC Dropout behaves like an inexpensive ensemble: each dropout mask creates a different “thinned” subnetwork. Aggregating many subnetworks typically improves calibration versus a single deterministic model, and the disagreement across subnetworks provides a useful confidence signal.

\textbf{Graduate level (Bayesian approximation).} Gal and Ghahramani showed that test-time dropout can be interpreted as approximate Bayesian inference, where each dropout mask corresponds to sampling from an approximate posterior distribution over network weights (a variational distribution implicitly defined by the dropout probability). This yields a practical estimate of \emph{epistemic uncertainty} (uncertainty due to limited data / model knowledge). \cite{gal2016dropout} In many applications, separating epistemic from \emph{aleatoric uncertainty} (inherent noise) is useful; Kendall and Gal provide a widely cited discussion of these uncertainty types and their role in deep learning systems. \cite{kendall2017uncertainties}

\textbf{How we use it for safety.} In the real-time loop, we run $N$ stochastic forward passes with dropout enabled, compute the mean class probability vector, and compute an uncertainty score from the sample distribution. A conservative policy is:
\begin{itemize}
\item require mean confidence above a threshold,
\item require uncertainty below a threshold, and
\item require stability across $k$ consecutive windows (hysteresis).
\end{itemize}
This reduces false positives and makes actuation safer, especially during transient artifacts (motion, blinks, poor electrode contact).

\section{PyTorch Primitives for Realtime Inference}
PyTorch primitives used in realtime inference:
\begin{itemize}
\item \texttt{\detokenize{torch.no_grad()}}
\item \texttt{\detokenize{model.eval()}} / \texttt{\detokenize{model.train()}}
\item \texttt{\detokenize{torch.softmax}}
\item batching
\item device placement
\item latency measurement
\end{itemize}

\todofig{Stability gate allow/deny timeline generated from Exercise 9 output.}{figs/stability_gate_timeline.pdf}

\section{Code Excerpts: MC Dropout and Safety Gating}

\begin{lstlisting}[language=PythonSafe,caption={Excerpt from \texttt{\detokenize{utils/mc_dropout.py}} L6--L32}]
def mc_dropout_predict(model, x, mc_samples=30):
    """
    Monte Carlo Dropout inference
    """

    model.train()  # IMPORTANT: enable dropout

    finger_probs = []
    action_probs = []

    with torch.no_grad():
        for _ in range(mc_samples):
            f_logits, a_logits = model(x)
            finger_probs.append(torch.softmax(f_logits, dim=1))
            action_probs.append(torch.softmax(a_logits, dim=1))

    model.eval()

    finger_probs = torch.stack(finger_probs)
    action_probs = torch.stack(action_probs)

    return {
        "finger_mean": finger_probs.mean(0),
        "finger_std": finger_probs.std(0),
        "action_mean": action_probs.mean(0),
        "action_std": action_probs.std(0),
    }
\end{lstlisting}

\begin{lstlisting}[language=PythonSafe,caption={Excerpt from \texttt{\detokenize{utils/inference.py}} L221--L258}]
        adaptive = min(
            0.99,
            max(
                self.config.base_threshold,
                self.config.base_threshold
                + self.config.uncertainty_weight * action_uncertainty,
            ),
        )

        self._stability.append(action_id)
        stability_ok = (
            len(self._stability) == self.config.stability_frames
            and len(set(self._stability)) == 1
        )

        velocity = 0.0
        if action_id != 0:
            velocity = action_confidence * (1.0 - action_uncertainty)

        prediction = {
            "action_id": action_id,
            "action_name": self.action_names.get(action_id, "UNKNOWN"),
            "finger_id": finger_id,
            "finger_name": self.finger_names.get(finger_id, "UNKNOWN"),
            "action_confidence": action_confidence,
            "action_uncertainty": action_uncertainty,
            "finger_confidence": finger_confidence,
            "finger_uncertainty": finger_uncertainty,
        }

        safety = {
            "base_threshold": self.config.base_threshold,
            "adaptive_threshold": adaptive,
            "allow_actuation": action_confidence >= adaptive and stability_ok,
            "stability_frames": self.config.stability_frames,
            "stability_ok": stability_ok,
            "velocity": velocity,
        }
\end{lstlisting}

\section{Exercises}

\exercise{9 (Coding)}{
Implement a stability gate: an action is emitted only if the top class is unchanged
for $k$ consecutive windows and mean confidence exceeds $\tau$.
Tune $(k,\tau)$ for your demo.

Run: \texttt{\detokenize{python exercises/ex09_stability_gate_demo.py --synthetic}}
}

% ============================================================

\chapter{Lesson 10: System Hardening, Testing, and Ethics}

\section{Hardening Checklist}

\begin{itemize}
\item reproducible environments and dependency pinning,
\item lint + format + unit tests,
\item timebase self-check as a runtime assertion,
\item cache validation and dataset identity enforcement,
\item safe-failure defaults (disable actuation on uncertainty).
\end{itemize}

\section{Ethics}

A prosthetic controller must preserve user agency:
\begin{itemize}
\item do not actuate confidently when uncertain,
\item provide transparent feedback when gated,
\item log failures for improvement.
\end{itemize}

\section{Exercises}

\exercise{10 (Reflection)}{
Write a one-page safety case: list hazards, mitigations, and verification steps for each mitigation.
}

% ============================================================

\appendix

\chapter{Solutions Manual}

\section{Exercise 1 (Conceptual)}
\solution{
\textbf{Explanation:} High accuracy can hide fundamental violations of the system contract.
\begin{itemize}
\item \textbf{Label-time misalignment:} labels are placed on the wrong samples due to clock drift/offset.
\item \textbf{Class imbalance deception:} REST dominates, so always predicting REST looks good but is useless.
\item \textbf{Evaluation leakage/cache mismatch:} cached predictions from a filtered dataset get reused on a different dataset.
\end{itemize}

\textbf{Expected output:} A list of at least three concrete failure mechanisms tied to timebase, label overlap, and evaluation integrity.

\textbf{Debugging notes:}
\begin{itemize}
\item Confirm timebase alignment with a self-check before training.
\item Inspect class balance, especially non-REST counts.
\item Verify cached prediction metadata against the current dataset filters.
\end{itemize}
}

\section{Exercise 2 (Conceptual)}
\solution{
\textbf{Explanation:} Movement intent has a temporal signature (preparation, execution, release) that is weak at any single instant.
A window integrates evidence and improves effective signal-to-noise ratio.

\textbf{Expected output:} A short explanation that mentions temporal structure and the benefit of integration over time.

\textbf{Debugging notes:}
\begin{itemize}
\item If your answer focuses only on spatial electrodes, revisit motor planning timing.
\item Tie the argument to real-time windowing instead of static snapshots.
\end{itemize}
}

\section{Exercise 3 (Coding)}
\solution{
\textbf{Explanation:} RMS energy is a coarse but fast proxy for high-amplitude artifacts.

\textbf{Expected output:}
\begin{itemize}
\item Standard output includes \texttt{\detokenize{windows=... flagged=...}} and an output JSON path.
\item In synthetic mode, at least one cluster of windows should be flagged.
\item Output file: \texttt{\detokenize{exercises/out/ex03_artifact_probe_intervals.json}}.
\end{itemize}

\textbf{Debugging notes:}
\begin{itemize}
\item If nothing is flagged, lower \texttt{\detokenize{--threshold}} or check \texttt{\detokenize{--fs}} and \texttt{\detokenize{--window-s}}.
\item If CSV parsing fails, confirm numeric columns and remove headers.
\end{itemize}
}

\section{Exercise 4 (Debugging)}
\solution{
\textbf{Explanation:} When events are shifted relative to samples, windows fail overlap tests and labels drift toward REST.

\textbf{Expected output:}
\begin{itemize}
\item \texttt{\detokenize{aligned_only}} should be greater than zero when the offset is large enough.
\item Output file: \texttt{\detokenize{exercises/out/ex04_timebase_break_and_measure.json}}.
\end{itemize}

\textbf{Debugging notes:}
\begin{itemize}
\item Ensure the offset is applied to events, not windows.
\item If counts do not change, increase \texttt{\detokenize{--offset-s}} or \texttt{\detokenize{--min-overlap-frac}}.
\end{itemize}
}

\section{Exercise 5 (Coding)}
\solution{
\textbf{Explanation:} Larger guard bands reduce ambiguous boundary windows but shrink the dataset.

\textbf{Expected output:}
\begin{itemize}
\item A table of guard/min-overlap settings with kept and dropped counts.
\item Output file: \texttt{\detokenize{exercises/out/ex05_guard_band_sweep.json}}.
\end{itemize}

\textbf{Debugging notes:}
\begin{itemize}
\item If everything drops, guard bands may exceed event length.
\item If everything keeps, minimum overlap may be too permissive.
\end{itemize}
}

\section{Exercise 6 (Debugging)}
\solution{
\textbf{Explanation:} Gap detection prevents windows with missing data from contaminating training.

\textbf{Expected output:}
\begin{itemize}
\item \texttt{\detokenize{gap_windows}} should be greater than zero when gaps are injected.
\item Output file: \texttt{\detokenize{exercises/out/ex06_gap_policy_sim.json}}.
\end{itemize}

\textbf{Debugging notes:}
\begin{itemize}
\item Ensure \texttt{\detokenize{--gap-s}} exceeds \texttt{\detokenize{--gap-threshold-s}}.
\item Increase \texttt{\detokenize{--gap-starts}} or \texttt{\detokenize{--duration-s}} if you do not see gaps.
\end{itemize}
}

\section{Exercise 7 (Conceptual)}
\solution{
\textbf{Explanation:} Joint action+finger supervision encourages a shared representation that encodes movement vs rest,
reducing spurious correlations that can dominate finger-only training.

\textbf{Expected output:} A brief explanation covering shared representation, regularization, and non-REST focus.

\textbf{Debugging notes:}
\begin{itemize}
\item If the answer ignores action labels, restate the multi-task benefit.
\item Mention evaluation on non-REST samples for finger accuracy.
\end{itemize}
}

\section{Exercise 8 (Coding)}
\solution{
\textbf{Explanation:} Adding a check for \texttt{\detokenize{timebase_version}} and \texttt{\detokenize{target_fs}} prevents reuse across incompatible windowing policies.

\textbf{Expected output:}
\begin{itemize}
\item Output includes \texttt{\detokenize{status=pass}} in synthetic mode and a JSON summary.
\item Output file: \texttt{\detokenize{exercises/out/ex08_cache_validator_extension.json}}.
\end{itemize}

\textbf{Debugging notes:}
\begin{itemize}
\item If JSON parsing fails, ensure \texttt{\detokenize{dataset_info}} is present in the npz.
\item When testing real data, pass \texttt{\detokenize{--dataset-info}} for the current dataset metadata.
\end{itemize}
}

\section{Exercise 9 (Coding)}
\solution{
\textbf{Explanation:} Stability gates add a discrete-time low-pass filter: only sustained, confident predictions trigger actuation.

\textbf{Expected output:}
\begin{itemize}
\item Output includes \texttt{\detokenize{allow_count=...}} and a JSON summary.
\item Output file: \texttt{\detokenize{exercises/out/ex09_stability_gate_demo.json}}.
\end{itemize}

\textbf{Debugging notes:}
\begin{itemize}
\item If nothing is allowed, lower \texttt{\detokenize{--confidence}} or \texttt{\detokenize{--stability-frames}}.
\item If too much is allowed, raise the threshold or stability frames.
\end{itemize}
}

\section{Exercise 10 (Reflection)}
\solution{
\textbf{Explanation:} A safety case links each hazard to a mitigation and a verification artifact.

\textbf{Expected output:} A one-page document that enumerates hazards, mitigations, and verification steps.

\textbf{Debugging notes:}
\begin{itemize}
\item Every mitigation should reference a concrete test, log, or checklist item.
\item Include both false trigger and missed trigger hazards.
\end{itemize}
}

\chapter{Figures}

\begin{itemize}
\item \textbf{Pipeline overview} (\texttt{\detokenize{figs/pipeline_overview.pdf}}): flow diagram from acquisition to real-time actuation.
\item \textbf{Somatotopy + montage} (\texttt{\detokenize{figs/montage_somatotopy.pdf}}): schematic of finger representation and electrode positions.
\item \textbf{Raw signal with events} (\texttt{\detokenize{figs/raw_with_events.pdf}}): EEG trace with blink artifact and movement interval overlay.
\item \textbf{Timebase alignment} (\texttt{\detokenize{figs/timebase_alignment.pdf}}): local clock and LSL clock with offset mapping.
\item \textbf{Window overlap policy} (\texttt{\detokenize{figs/window_overlap_policy.pdf}}): window, event, overlap fraction, and guard band geometry.
\item \textbf{Resampling grid} (\texttt{\detokenize{figs/resampling_grid.pdf}}): irregular timestamps mapped to uniform grid.
\item \textbf{Timebase misalignment impact} (\texttt{\detokenize{figs/timebase_alignment_effect.pdf}}): bar chart aligned vs offset window retention from Exercise 4.
\item \textbf{Model architecture} (\texttt{\detokenize{figs/model_architecture.pdf}}): shared backbone with action and finger heads.
\item \textbf{Reliability diagram} (\texttt{\detokenize{figs/reliability_diagram.pdf}}): confidence bins vs accuracy with ECE.
\item \textbf{Latency vs accuracy} (\texttt{\detokenize{figs/latency_accuracy_tradeoff.pdf}}): tradeoff curve with chosen operating point.
\item \textbf{Stability gate timeline} (\texttt{\detokenize{figs/stability_gate_timeline.pdf}}): allow/deny trace over frames from Exercise 9.
\end{itemize}

\chapter{Reference Links}

\begin{itemize}
\item Lab Streaming Layer docs/repo for clock functions and timing model. \url{https://labstreaminglayer.org/} \url{https://github.com/sccn/labstreaminglayer}
\item ERD/ERS foundational review for mu/beta movement modulation. \url{https://doi.org/10.1016/S1388-2457(99)00185-5}
\item Calibration + reliability diagrams (ECE). \url{https://arxiv.org/abs/1706.04599}
\item MC Dropout uncertainty. \url{https://arxiv.org/abs/1506.02142}
\item Consumer EEG constraints + sampling rate context (Muse-class). \url{https://doi.org/10.3389/fnins.2017.00109}
\item PyTorch documentation for model building and training utilities used in the pipeline. \url{https://docs.pytorch.org/docs/stable/index.html}
\end{itemize}

\chapter*{Final Reflection}

\begin{quote}
Correctness, discipline, and humility matter more than model complexity.
\end{quote}

\end{document}


\section{References}
\begin{thebibliography}{9}

\bibitem{LSL}
Christian A. Kothe.
\newblock \emph{Lab Streaming Layer (LSL)}.
\newblock Project documentation and source: \url{https://github.com/sccn/labstreaminglayer}.

\bibitem{pylsl}
\emph{pylsl: Python bindings for Lab Streaming Layer}.
\newblock Source and documentation: \url{https://github.com/labstreaminglayer/pylsl}.

\bibitem{Bleak}
\emph{bleak: A cross-platform Python BLE library}.
\newblock Source and documentation: \url{https://github.com/hbldh/bleak}.

\bibitem{Muse2spec}
\emph{Muse 2 EEG technical notes (channels and sampling rate)}.
\newblock Summary of EEG channels and 256\,Hz sampling commonly used in Muse 2 tooling: \url{https://choosemuse.com/muse-2-specifications/}.


\end{thebibliography}

